% !TeX root = ../Cosmology.tex

%% --------------------------------------------------------
\chapter{Кінематика інфляції}
\label{sec:inflation_kinematics}
%% --------------------------------------------------------

У попередньому розділі ми сформулювали три фундаментальні проблеми
стандартної космології: плоскості, горизонту і первинних збурень.
Цей розділ показує, як усі три вирішуються \emph{кінематично}~---
без залучення конкретної мікрофізичної моделі~---
якщо припустити, що до гарячої стадії Великого вибуху існувала
коротка епоха прискореного розширення: \textbf{інфляція}.
Мікрофізику~--- що саме рухає інфляцію~--- буде розглянуто в наступному
розділі.

\medskip\noindent
\textbf{Логічна карта розділу}
\begin{center}
	\begin{tikzpicture}[roadmap]
		\node[box] (prob) {3 проблеми $\Lambda$CDM};
		\node[box, right=of prob] (kin) {Інфляція: $\ddot{a} > 0$};
		\node[box, above right=0.3cm and 0.8cm of kin] (hor) {Горизонт};
		\node[box, below right=0.3cm and 0.8cm of kin] (flat) {Плоскість};
		\node[box, right=4.5cm of kin, fill=red!5, draw=red!70!black]
		(efolds) {$N \gtrsim 60$};
		\draw[arrow] (prob) -- (kin);
		\draw[arrow] (kin) -- (hor);
		\draw[arrow] (kin) -- (flat);
		\draw[arrow] (hor) -- (efolds);
		\draw[arrow] (flat) -- (efolds);
		\node[font=\tiny\bfseries, text=blue!70!black, above=0.1cm of prob, align=center]
		{плоскість, горизонт, збурення};
		\node[font=\tiny\bfseries, text=blue!70!black, above=0.1cm of kin, align=center]
		{конформний час $\eta < 0$};
		\node[font=\scriptsize, gray, below=0.1cm of efolds] {необхідна тривалість};
	\end{tikzpicture}
\end{center}

%% --------------------------------------------------------
\section{Конформний час під час інфляції і горизонт Хаббла}
\label{sec:conformal_inflation}
%% --------------------------------------------------------

Ключ до розуміння інфляції~--- конформний час $\eta$.
Нагадаємо (§\,\ref{sec:horizon}): під час інфляції
$a(t) = a_i e^{H_{\text{Inf}} t}$ при $H_{\text{Inf}} \approx \mathrm{const}$,
тому:
\begin{equation}\label{eq:eta_inflation}
	\eta(t) = \int_{-\infty}^{t}\frac{dt'}{a_i e^{H_{\text{Inf}} t'}}
	= -\frac{1}{a(t)H_{\text{Inf}}} = -\frac{1}{\mathcal{H}}.
\end{equation}
Вся інфляційна епоха відповідає $\eta \in (-\infty, 0)$,
а конформний горизонт Хаббла:
\begin{equation}\label{eq:conformal_horizon}
	\mathcal{H}^{-1} = \frac{1}{a H_{\text{Inf}}} = -\eta
\end{equation}
\emph{спадає} з часом~--- саме це є кінематичним визначенням інфляції.
Фізично: масштаби виходять за горизонт швидше, ніж встигають
обмінятись сигналами.

На рисунку~\ref{fig:conformal_inflation} та в подальших викладках
ми позначаємо $\eta_{\text{Init}}$ конформний час у момент,
коли найбільший спостережуваний сьогодні масштаб перетнув горизонт.
Це скінченне число, яке і визначає $\Delta\eta_{\text{Inf}}$.

Введемо позначення для ключових моментів конформного часу (рис.~\ref{fig:conformal_inflation}):
\begin{align*}
	\eta_\text{Init}   & \to -\infty                    \\
	\eta_\text{End}    & = 0 \quad \text{--- кінець інфляції} \\
	\eta_r             & > 0 \quad \text{--- розігрів,}                                                         \\
	\eta_0             & > \eta_r \quad \text{--- сьогодні.}
\end{align*}
Конформні інтервали двох епох:
\begin{equation}\label{eq:conformal_intervals}
	\Delta\eta_{\text{Inf}} \equiv \eta_{\text{End}} - \eta_{\text{Init}}, \qquad
	\Delta\eta_{\text{Bb}}  \equiv \eta_0 - \eta_{\text{End}}.
\end{equation}

%---------------------------------------------------------
\begin{figure}[h!]
\centering
\begin{tikzpicture}[
    point/.style      = {circle, fill=red, inner sep=1.5pt},
    label_text/.style = {font=\scriptsize, align=center},
    arrow_line/.style = {-{Stealth}, thick},
]

% =============================================================
% ПАРАМЕТРИ — міняй тільки тут
% =============================================================
\def\etaInit{-6}   % η_Init (відносно η_End ≡ 0)
\def\etaEnd {0}      % η_End  ≡ 0 (вісь y)
\def\etaR   {1.8}    % η_r  (рекомбінація)
\def\etaZero{5.0}    % η_0  (сьогодні)
\def\yMax   {4.2}    % верхня межа осі y
\def\axisExt{0.5}    % виступ осей за межі

% Нормування кривих (підбираються один раз)
\def\scaleA {0.55}   % нахил a∝η після інфляції
\def\scaleH {0.90}   % нахил H⁻¹∝η після інфляції
\def\scaleAi{0.55}   % нормування гіперболи a під час інфляції

% =============================================================
% ОСІ
% =============================================================
\draw[arrow_line]
    ({\etaInit - \axisExt}, 0) -- ({\etaZero + \axisExt}, 0)
    node[right] {$\eta$};

\draw[arrow_line]
    (0, -0.3) -- (0, {\yMax + \axisExt})
    node[above, font=\small] {};

% Підписи функцій біля осі y
\node[theoryred, font=\small, left=4pt] at (0, \yMax) {$a(\eta)$};
\node[accent,    font=\small, right=4pt] at (0, \yMax) {$\mathcal{H}^{-1}(\eta)$};

% =============================================================
% КЛЮЧОВІ ВУЗЛИ НА ОСІ η
% =============================================================
\node[point, label=above:{$\eta_{\text{Init}}$}] (nInit) at (\etaInit, 0) {};
\node[point, label=above:{$\eta_{\text{End}}$}] (nEnd) at (\etaEnd, 0) {};
\node[point, label=above:{$\eta_r$}]  (nR)    at (\etaR,    0) {};
\node[point, label=above:{$\eta_0$}]  (nZero) at (\etaZero, 0) {};

% =============================================================
% КРИВІ — ІНФЛЯЦІЙНА ЕПОХА  η ∈ [η_Init, η_End)
%   a(η)        = scaleAi / (0.15 - η)   → гіпербола, ↑ при η→0⁻
%   H⁻¹(η)     = scaleH  * (-η)          → лінійна,  ↓ при η→0⁻
% =============================================================
\draw[theoryred, thick, domain={\etaInit}:{-0.05}, samples=120]
    plot (\x, {\scaleAi / (0.15 - \x)});

\draw[accent, thick, domain={\etaInit}:{-0.05}, samples=80]
    plot (\x, {\scaleH * (-\x)});

% =============================================================
% КРИВІ — ПОСТ-ІНФЛЯЦІЙНА ЕПОХА  η ∈ (η_End, η_0]
%   a(η)    ∝ η  (радіаційна)
%   H⁻¹(η) ∝ η  (зростає, різний нахил)
% =============================================================
\draw[theoryred, thick, domain={0.05}:{\etaZero}, samples=80]
    plot (\x, {\scaleA * \x});

\draw[accent, thick, domain={0.05}:{\etaZero}, samples=80]
    plot (\x, {\scaleH * \x});

% =============================================================
% КОНФОРМНІ ІНТЕРВАЛИ (рівень \yMax - 0.3)
% =============================================================
\pgfmathsetmacro{\arrowY}{- 0.6}

\draw[<->, thick, green!50!black]
    ({\etaInit + 0.1}, \arrowY) -- ({\etaEnd - 0.1}, \arrowY);
\node[font=\footnotesize, green!50!black, above] at ({(\etaInit)/2}, \arrowY)
    {$\Delta\eta_{\text{Inf}}$};

\draw[<->, thick, green!50!black]
    ({\etaEnd + 0.1}, \arrowY) -- ({\etaZero - 0.1}, \arrowY);
\node[font=\footnotesize, green!50!black, above] at ({(\etaR + \etaZero)/2}, \arrowY)
    {$\Delta\eta_{\text{Bb}}$};

% =============================================================
% ЛЕГЕНДА (прив'язана до nEnd + зсув)
% =============================================================
\pgfmathsetmacro{\legX}{\etaR + 2}
\pgfmathsetmacro{\legYa}{\yMax - 0.5}
\pgfmathsetmacro{\legYb}{\yMax - 0.85}

\draw[theoryred, thick] (\legX, \legYa) -- ({\legX + 0.8}, \legYa)
    node[right, black, font=\scriptsize] {$a(\eta)$};
\draw[accent,    thick] (\legX, \legYb) -- ({\legX + 0.8}, \legYb)
    node[right, black, font=\scriptsize] {$\mathcal{H}^{-1}(\eta)$};

% =============================================================
% ПІДПИСИ ПОДІЙ ЗНИЗУ
% =============================================================
\node[label_text, below=1.5cm of nInit] (descInit)
    {Початок інфляції\\(механізм~--- відкрите питання)};
\draw[arrow_line] (descInit.north) -- (nInit);

\node[label_text, below=1cm of nEnd] (descEnd)
    {кінець інфляції,\\гарячий Великий вибух};
\draw[arrow_line] (descEnd.north) -- (nEnd);

\node[label_text, below=0.6cm of nR]    (descR)    {Рекомбінація};
\draw[arrow_line] (descR.north) -- (nR);

\node[label_text, below=1cm of nZero] (descZero) {Наш час};
\draw[arrow_line] (descZero.north) -- (nZero);

% =============================================================
% РАМКА: умова вирішення проблеми горизонту
% =============================================================
%\node[draw=red!70!black, rounded corners=3pt, fill=red!5,
%      font=\small\bfseries, text=red!70!black, align=center,
%      anchor=north west] at ({\etaInit}, {\yMax + \axisExt - 0.05})
%    {$\Delta\eta_{\text{Inf}} \gg \Delta\eta_{\text{Bb}}$\\
%     \footnotesize вирішення проблеми горизонту};

\end{tikzpicture}
\caption{Еволюція масштабного фактора $a(\eta)$ (червона)
    та конформного горизонту Хаббла $\mathcal{H}^{-1} = (aH)^{-1}$
    (синя) у конформному часі. Вісь~$y$ проходить через
    $\eta_{\text{End}}$. Під час інфляції $a$ зростає,
    але $\mathcal{H}^{-1}$ \emph{спадає}~--- масштаби виходять
    за горизонт і «замерзають». Після інфляції обидві величини
    зростають і заморожені масштаби повертаються в горизонт.}
\label{fig:conformal_inflation}
\end{figure}
%---------------------------------------------------------

%% --------------------------------------------------------
\section{Проблема горизонту і умова її вирішення}
\label{sec:horizon_problem}
%% --------------------------------------------------------

\textbf{Нагадування проблеми.}
CMB є однорідним з точністю $10^{-5}$ по всьому небу.
Але в стандартній гарячій космології без інфляції ділянки неба,
розділені більше ніж $\sim 2^\circ$, ніколи не перебували
у причинному контакті: їхні горизонти частинок на момент
рекомбінації не перетиналися. Звідки тоді однакова температура?

\textbf{Розв'язання через конформний час.}
Горизонт частинок~--- це максимальна відстань, з якої
міг прийти сигнал від початку Всесвіту до моменту $\eta$:
\begin{equation*}
	l_H(\eta) = a(\eta)\int_{\eta_{\text{Init}}}^{\eta} d\eta'
	= a(\eta)\,(\eta - \eta_{\text{Init}}).
\end{equation*}
Розмір горизонту на кінці інфляції $\eta_{\text{End}}$,
\emph{розтягнутий} до сьогоднішнього розміру:
\begin{equation}\label{eq:horizon_today}
	l_{H,\text{End}}(\eta_0) = a_0\,(\eta_{\text{End}} - \eta_{\text{Init}})
	= a_0\,\Delta\eta_{\text{Inf}}.
\end{equation}
Розмір нашої спостережуваної Всесвіту сьогодні:
\begin{equation}\label{eq:observable_universe}
	l_{H,0} = a_0\,(\eta_0 - \eta_{\text{End}}) = a_0\,\Delta\eta_{\text{Bb}}.
\end{equation}
Щоб весь спостережуваний Всесвіт поміщався всередині одного
причинно-зв'я\-за\-ного регіону з часів інфляції, потрібно:
\begin{equation}\label{eq:horizon_condition}
	\boxed{l_{H,\text{End}}(\eta_0) \gg l_{H,0}
		\quad\Leftrightarrow\quad
		\Delta\eta_{\text{Inf}} \gg \Delta\eta_{\text{Bb}}.}
\end{equation}
Тобто конформний інтервал інфляції має бути набагато більшим
за весь пост-інфляційний конформний час.

\textbf{Чи виконується умова~\eqref{eq:horizon_condition}?}
Обчислимо відношення. Для експоненціальної інфляції $a(t) = a_i e^{H_{\text{Inf}} t}$:
\begin{equation*}
	\Delta\eta_{\text{Inf}}
	= \int_{t_{\text{Init}}}^{t_{\text{End}}} \frac{dt}{a_i e^{H_{\text{Inf}} t}}
	= \frac{1}{H_{\text{Inf}} a_i}\left(e^{-H_{\text{Inf}}t_{\text{Init}}}
	  - e^{-H_{\text{Inf}}t_{\text{End}}}\right)
	\approx \frac{1}{H_{\text{Inf}} a_{\text{Init}}},
\end{equation*}
де останній крок використовує те, що підінтегральна функція
$e^{-H_{\text{Inf}}t}$ \emph{спадає}: внесок нижньої межі
$e^{-H_{\text{Inf}}t_{\text{Init}}}$ переважає над верхньою
$e^{-H_{\text{Inf}}t_{\text{End}}} = e^{-N} e^{-H_{\text{Inf}}t_{\text{Init}}} \ll e^{-H_{\text{Inf}}t_{\text{Init}}}$.
Пост-інфляційний інтервал оцінюється як $\Delta\eta_{\text{Bb}} \sim 1/(a_0 H_0)$. Тоді:
\begin{equation}\label{eq:horizon_ratio}
	\frac{\Delta\eta_{\text{Inf}}}{\Delta\eta_{\text{Bb}}} \sim e^N \cdot \frac{a_0 H_0}{a_{\text{End}} H_{\text{Inf}}} \gg 1.
\end{equation}

Щоб визначити необхідну тривалість інфляції ($N \equiv \ln(a_{\text{End}}/a_{\text{Init}})$),
врахуємо зв'язок масштабів через температуру.
Використовуючи закон червоного зміщення ($a \sim 1/T$),
отримуємо $\frac{a_{\text{End}}}{a_0} \sim \frac{T_0}{T_{\text{reh}}}$.
Підставляючи це в умову $\frac{\Delta\eta_{\text{Inf}}}{\Delta\eta_{\text{Bb}}} \gg 1$, маємо:
\begin{equation*}
	\frac{a_{\text{End}}}{a_{\text{Init}}} \gg \frac{T_0}{H_0} \cdot \frac{H_{\text{Inf}}}{T_{\text{reh}}}.
\end{equation*}
Логарифмуючи цей вираз, знаходимо нижню межу для кількості $e$-foldів:
\begin{equation*}
	N \gg \ln \left( \frac{T_0}{H_0} \right) + \ln \left( \frac{H_{\text{Inf}}}{T_{\text{reh}}} \right).
\end{equation*}

Щоб отримати числову оцінку, потрібно знати $H_{\text{Inf}}$.
За порядком величини природно вважати, що максимально можливий
масштаб інфляції --- планківський: $H_{\text{Inf}} \lesssim M_{\text{Pl}}$.
Це не сувора межа, а розумна верхня оцінка: якби $H_{\text{Inf}}$ було
набагато меншим, то й $N$ було б більшим, що лише посилює виконання
умов. Тому для оцінки мінімально необхідного $N$ ми можемо взяти
$H_{\text{Inf}} \sim M_{\text{Pl}}$.

Тоді:
\begin{equation*}
	N \gg 67 + \ln\left(\frac{M_{\text{Pl}}}{T_{\text{reh}}}\right).
\end{equation*}

Для типових значень $T_{\text{reh}} \sim 10^{15}$~GeV маємо
$\ln(M_{\text{Pl}}/T_{\text{reh}}) \approx -4.6$, звідки $N \gg 62.4$,
тобто $N \gtrsim 60$.

\begin{keybox}{Залежність від температури розігріву}
	Оцінка $N$ чутливо залежить від $T_{\text{reh}}$:
	\begin{itemize}
		\item Високий розігрів ($T_{\text{reh}} \sim 10^{15}$~GeV):
		      $\ln(M_{\text{Pl}}/T_{\text{reh}}) \approx -4.6$,
		      $N \gtrsim 62$.
		\item Типовий діапазон ($T_{\text{reh}} \sim 10^{13}$~GeV):
		      $\ln(M_{\text{Pl}}/T_{\text{reh}}) \approx -7.8$,
		      $N \gtrsim 59$.
		\item Низький розігрів ($T_{\text{reh}} \sim 10^3$~GeV):
		      $\ln(M_{\text{Pl}}/T_{\text{reh}}) \approx -34.5$,
		      $N \gtrsim 33$.
	\end{itemize}
	Оскільки більшість реалістичних моделей тяжіють до високих
	масштабів енергій ($10^{13}$--$10^{15}$~GeV),
	у космології як стандартне базове значення приймають
	$N \gtrsim 55$--$65$.
\end{keybox}

Таким чином, у природних сценаріях інфляції вимога $N \gtrsim 60$
гарантовано вирішує проблему горизонту.

%% --------------------------------------------------------
\section{Проблема плоскості і мінімальна кількість $e$-foldів}
\label{sec:flatness_efolds}
%% --------------------------------------------------------

З розділу~\ref{subsec:flatness_paradox} нагадаємо, що під час інфляції
($w \approx -1$) відхилення від плоскості \emph{спадає}:
\begin{equation*}
	|\Omega_{\text{tot}} - 1| = \frac{k}{a^2 H^2} \propto e^{-2H_{\text{Inf}} t}.
\end{equation*}
Після $N$ $e$-foldів початкове відхилення $|\Omega - 1|_i \sim 1$
пригнічується до:
\begin{equation*}
	|\Omega_{\text{tot}} - 1|_{\text{End}} \approx e^{-2N}.
\end{equation*}
Щоб сьогодні мати $|\Omega_k| < 0.005$, враховуючи, що після інфляції
відхилення \emph{зростає} як $a^2$ (§\,\ref{subsec:flatness_paradox}),
отримуємо умову:
\begin{equation}\label{eq:flatness_efolds}
	e^{-2N} \cdot \left(\frac{a_0}{a_{\text{End}}}\right)^2 \lesssim 0.005.
\end{equation}
Підставляючи $a_0/a_{\text{End}} = T_{\text{reh}}/T_0$
(із закону збереження ентропії $a \propto 1/T$) і
$T_{\text{reh}} \sim 10^{15}$~GeV, $T_0 \sim 10^{-4}$~eV, знаходимо:
\begin{equation*}
	\frac{a_0}{a_{\text{End}}} \sim \frac{10^{15}\,\text{GeV}}{10^{-13}\,\text{GeV}}
	= 10^{28} \approx e^{64},
\end{equation*}
тому умова~\eqref{eq:flatness_efolds} дає:
\begin{equation}\label{eq:N_min}
	\boxed{N \gtrsim 60.}
\end{equation}

%% --------------------------------------------------------
\section{Проблема первинних збурень}
\label{sec:primordial_perturbations_bridge}
%% --------------------------------------------------------

\textbf{Нагадування проблеми.}
Модель $\Lambda$CDM приймає початковий спектр флуктуацій густини
як феноменологічний вхідний параметр ($A_s \approx 2.1\times10^{-9}$,
$n_s \approx 0.965$). Але \enquote{звідки} взялися ці флуктуації
і чому спектр майже степеневий~--- стандартна космологія мовчить.

\textbf{Роль інфляції.}
Кінематика інфляції створює необхідні умови для генерації збурень:
експоненціальне розтягнення простору перетворює мікроскопічні квантові
флуктуації в макроскопічні неоднорідності. Однак сам механізм
народження збурень потребує залучення мікрофізики~---
квантової теорії поля інфлатона. Це буде детально розглянуто
в розділі~\ref{sec:quantum}. Тут лише зазначимо, що інфляція
не лише вирішує проблеми горизонту і плоскості, а й забезпечує
природне походження первинних неоднорідностей~--- головне
джерело великомасштабної структури Всесвіту.

%% --------------------------------------------------------
\section{Місток: від кінематики до мікрофізики}
\label{sec:kinematics_bridge}
%% --------------------------------------------------------

Ми встановили, що інфляція \emph{кінематично} вирішує проблеми
горизонту і плоскості за умови $N \gtrsim 60$, а також створює
умови для генерації первинних збурень.
Але залишається відкритим питання: \emph{що} фізично рухає інфляцію?

Прискорене розширення $\ddot{a} > 0$ вимагає (з другого рівняння
Фрідмана~\eqref{eq:F2}):
\begin{equation*}
	\rho + \frac{3p}{c^2} < 0 \quad\Rightarrow\quad w < -\frac{1}{3}.
\end{equation*}
Жодна відома форма звичайної матерії не задовольняє цю умову.
Тому вводять \textbf{інфлатон}~--- скалярне поле $\phi$ з потенціалом
$V(\phi)$, яке при повільному скоченні вздовж пологого схилу
забезпечує $w \approx -1$.
Динаміці інфлатона присвячений наступний розділ.

			Інфляція --- це кінематичний механізм, який при $N \gtrsim 60$:
			\begin{enumerate}
				\item розв'язує проблему горизонту ($\Delta\eta_{\text{Inf}} \gg \Delta\eta_{\text{Bb}}$),
				\item розв'язує проблему плоскості (експоненціальне пригнічення $\Omega_k$),
				\item створює умови для генерації первинних збурень.
			\end{enumerate}
			Мікрофізична реалізація вимагає скалярного поля --- інфлатона.

%% --------------------------------------------------------
\section*{Задачі}
\addcontentsline{toc}{section}{Задачі}
%% --------------------------------------------------------

\begin{problem}{Конформний час під час інфляції}
	Показати, що для де-Сіттерівського розширення
	$a(\eta) = -1/(H_{\text{Inf}}\eta)$ при $\eta \in (-\infty, 0)$.
	Перевірити, що $\mathcal{H} = a H_{\text{Inf}} = -1/\eta$.

	\textit{Вказівка:} підставте $a(t) = a_i e^{H_{\text{Inf}}t}$
	у визначення конформного часу $d\eta = dt/a(t)$ і проінтегруйте.
	Нормування обрати так, щоб $\eta \to 0^-$ при $t \to +\infty$.

	\textit{Відповідь:} $\eta = -\frac{1}{a(t)H_{\text{Inf}}}$,
	звідки $a(\eta) = -\frac{1}{H_{\text{Inf}}\eta}$
	і $\mathcal{H} \equiv \frac{a'}{a} = -\frac{1}{\eta} = aH_{\text{Inf}}$.
\end{problem}

% -------------------------------------------------------------------

\begin{problem}{Мінімальна кількість $e$-foldів}
	Оцінити мінімальне $N$ з умови, що горизонт частинок на кінці
	інфляції, розтягнутий до сьогодні, перевищує розмір
	спостережуваного Всесвіту $c/H_0$.
	Використати $T_{\text{rh}} \sim 10^{15}$~GeV і $T_0 \sim 10^{-4}$~eV.
\end{problem}

% -------------------------------------------------------------------

\begin{problem}{Проблема горизонту без інфляції}
	Оцінити кут на небі, що відповідає горизонту частинок на момент
	рекомбінації ($z_* \approx 1100$) у стандартній космології без інфляції.
	Порівняти з кутовим розміром однорідної ділянки CMB.

	\textit{Вказівка:} фізичний розмір горизонту на рекомбінації
	$l_H(z_*) \approx c/H(z_*)$, де $H(z_*) \approx H_0\,\Omega_m^{1/2}(1+z_*)^{3/2}$.
	Кутовий розмір~--- це $\theta \approx l_H(z_*)/d_A(z_*)$,
	де $d_A(z_*) \approx c/H_0 \cdot 2/(1+z_*)^{1/2}$.

	\textit{Відповідь:} $\theta \approx \Omega_m^{1/2}/(1+z_*)^{1/2} \approx 1.7^\circ$~---
	саме стільки охоплює один причинно-зв'язаний патч на карті CMB.
	Оскільки таких патчів $\sim (180^\circ/1.7^\circ)^2 \sim 10^4$,
	виникає питання: чому всі вони мають однакову температуру?
\end{problem}

% -------------------------------------------------------------------

\begin{problem}{Об'єднана умова на $N$ з плоскості}
	Використовуючи вирази для еволюції $|\Omega_{\text{tot}}-1|$ під час
	інфляції ($\propto e^{-2N}$) та після ($\propto a^2$),
	отримайте умову на $N$ з вимоги $|\Omega_k(t_0)| < 0.005$.
	Оцініть мінімальне $N$ для $T_{\text{reh}} = 10^{15}$~GeV та $T_{\text{reh}} = 10^3$~GeV.
	\textit{Відповідь:} $N \gtrsim 60$ для високого розігріву,
	$N \gtrsim 30$ для низького.
\end{problem}